\chapter{Analytische Geometrische}

\section{Resourecen}

\subsection{Kerncurriculum}

https://cuvo.nibis.de/cuvo.php?p=download&upload=208

\startitemize[packed]
\item Raumanschauung und Koordinatisierung
\startitemize[packed]
  \item Punkte und Vektoren in Ebene und Raum durch Tupel beschreiben
  \item die bildliche Darstellung und Koordinatisierung zur Beschreibung von Punkten, Strecken, ebenen Flächen und einfachen Körpern nutzen
  \item Addition, Subtraktion und skalare Multiplikation von Vektoren anwenden und geometrisch veranschaulichen
  \item Kollinearität zweier Vektoren überprüfen
\stopitemize
\item Darstellungsformen
\startitemize[packed]
  \item Geraden- und Ebenengleichungen in Parameterform verwenden
  \item Ebenengleichungen in Normalen- und Koordinatenform verwenden
  \item zwischen den Darstellungsformen wechseln
\stopitemize
\item Maße und Lagen
\startitemize[packed]
  \item Abstände zwischen Punkten, Geraden und Ebenen bestimmen
  \item Skalarprodukt geometrisch als Ergebnis einer Projektion deuten und verwenden
  \item Orthogonalität zweier Vektoren überprüfen
  \item Winkelgrößen bestimmen
  \item Lagebeziehungen von Geraden, Geraden und Ebenen sowie von Ebenen untersuchen und Schnittprobleme lösen
  \item den Gauß-Algorithmus zur Lösung linearer Gleichungssysteme erläutern und in geeigneten Fällen anwenden
\stopitemize
\stopitemize

\section{Punkte und Vektoren}

\startbuffer[analytischegeometrie:punkte_vektoren:darstellung_vektor_punkt]
\startMPpage
	u := 1cm;
	drawarrow (0, -0.5u)--(0, 4u);
	drawarrow (-0.5u, 0)--(4u, 0);

	label.lft(btex $x_2$ etex, (0, 4u));
	label.bot(btex $x_1$ etex, (4u, 0));

	drawarrow (0, 0)--(1u, 3u);
	label.rt(btex $\vec{a}=\pmatrix{1;3}$ etex, (1u, 3u));

	drawdot (3u, 1u) withpen pencircle scaled 3pt;
	label.bot(btex $P(3|1)$ etex, (3u, 1u));
\stopMPpage
\stopbuffer

\placefigure[left][analytischegeometrie:punkte_vektoren:darstellung_vektor_punkt]{Darstellung Vektor und Punkt}{
\typesetbuffer[analytischegeometrie:punkte_vektoren:darstellung_vektor_punkt]
}

Ein Punkt beschreibt einen Ort in einem Raum. Dafür beeinhaltet er Koordinaten für alle Achsen:

\startformula
P(x_1|x_2|x_3)
\stopformula

Ein Vektor hingegen {\bf zeigt} in eine bestimmte Richtung und hat dabei ein feste Länge. Wie ein Punkt hat eine Vektor angaben

\startformula
\vec{a}=\pmatrix{x_1;x_2}
\stopformula

\startbuffer[analytischegeometrie:punkte_vektoren:satzdespythagora]
\startMPpage

u := 1cm;
pair A, B, C;
A:=(0, 0);
B:=(0, 2u);
C:=(2u, 0);

draw A--B--C--cycle;
\stopMPpage
\stopbuffer

\margintext{
\placefigure[here][analytischegeometrie:punkte_vektoren:satzdespythagora]{Satz des Pythagoras}{
\typesetbuffer[analytischegeometrie:punkte_vektoren:satzdespythagora]
}}

Die Länge eines Vektors berechnet man, in dem man den Satz des
Pythagoras\sidenote{HI} mehrmals anwendet. So gesehen wendet
man ihn erst auf den $x_1-x_2$-Abschnitt an und danach auf den
$x_2-x_3$-Abschnitt.

\startplaceformula[analytischegeometrie:punkte_vektoren:laenge_vektor]
\startformula
|\vec{a}|=\sqrt{x_1^2+x_2^2+x_3^2}
\stopformula
\stopplaceformula

Das gleiche gilt auch für den Abstand zwischen zwei Punkten. Dafür muss man allerding erst den Vektor zwischen den beiden Punkten ermitteln.


